\section{Introducción}
\label{sec:introduccion}

\subsection{Presentación general}
\label{subsec:intro-presentacion}

Este trabajo implementa un videojuego de plataformas en 2D. El núcleo del sistema es una simulación
discreta por \emph{frame}: en cada paso se recibe un \emph{delta time} y una entrada externa,
y se produce un nuevo \emph{world}, instantánea del estado en el \emph{game plane} con el
\emph{player}, los \emph{enemy}, las plataformas y los ítems aún presentes en el nivel.
El contenido estático de cada etapa proviene de una \emph{level definition} cargada antes de
jugar. Durante la partida el motor actualiza ese \emph{world} frame a frame.

La interacción con el \emph{medio} ocurre en la frontera del
sistema: entrada del usuario, lectura de datos externos y rendering en tiempo real. Esas
responsabilidades quedan fuera del núcleo de simulación. La lógica de juego se expresa como
transformaciones puras sobre el \emph{world}, mientras que la orchestration de estado global,
configuración y efectos se delega a capas superiores del motor.

\subsection{Motivación}
\label{subsec:intro-motivacion}

El objetivo central del trabajo es demostrar, en un dominio no trivial, los conceptos
vistos en la materia: funciones puras, composición con tipos algebraicos, abstracción y
manejo de efectos con mónadas. Un juego de plataformas resulta un caso de estudio exigente
porque combina estado global, actualización continua y entrada en tiempo real.

Elegimos un videojuego también por una razón práctica: el feedback visual hace
tangible el comportamiento del programa. Cada \emph{frame} se puede observar de inmediato,
y un error en física, colisiones o comportamiento de un \emph{enemy} se manifiesta en pantalla
sin depender de trazas ni dumps de estado. Eso acorta el ciclo de validación y mantiene el
foco en la corrección del modelo, no solo en que el código compile.

\subsection{Objetivos}
\label{subsec:intro-objetivos}

\textbf{Objetivo general.}
Demostrar conceptos de Programación Funcional mediante la construcción de un videojuego de
plataformas~2D jugable, con un motor modular que separa simulación pura, orchestration con
mónadas e integración con el entorno externo.

\textbf{Objetivos específicos.}

\begin{enumerate}
  \item Implementar un juego~2D jugable con las mecánicas centrales del dominio: movimiento
    del \emph{player} con gravedad y salto, apoyo sobre plataformas estáticas y móviles,
    recolección de ítems con efecto sobre puntaje, y combate básico con ataques
    cuerpo a cuerpo y a distancia. El \emph{run} debe soportar pérdida de vidas,
    \emph{respawn} en un \emph{spawn point} y condición de victoria \emph{híbrida}: llegar a la
    \emph{exit zone} con el puntaje mínimo exigido y, si hay jefe en el nivel, derrotarlo.

  \item Diseñar un motor modular por capas (\texttt{Domain}, \texttt{UseCases},
    \texttt{Adapters}, \texttt{Frameworks}), de modo que la simulación viva en un núcleo
    puro sin \texttt{IO}, la orchestration del \emph{frame} quede en casos de uso, y la
    entrada, el reloj y el dibujado se concentren en adaptadores concretos. La transición
    central del motor se expresa como una función pura de paso sobre el \emph{world}.

  \item Modelar los conceptos del dominio del juego en tipos algebraicos y \emph{value objects}
    dentro de \texttt{Domain/}: distinguir \emph{level definition} y \emph{world}, nombrar
    fases de partida, magnitudes físicas y \emph{behaviour tuning} (\texttt{BehaviourTuning}) como
    valores con invariantes, en lugar de repartir reglas en booleanos o strings sueltos.

  \item Modelar comportamientos de \emph{entity} (en particular \emph{enemy}) mediante un
    \emph{DSL} basado en \emph{free monad}: el \emph{AST} describe
    acciones como desplazarse, esperar o atacar, y un interpreter puro en el dominio las ejecuta
    sobre el \emph{world}. Esa separación permite combinar patrones de IA (patrol,
    chase, disparo a distancia) sin mezclar lógica de comportamiento con física ni con
    efectos del entorno.

  \item Orquestar el estado global del juego, la configuración de física y parámetros de
    partida, y los errores recuperables al cargar datos, con una pila monádica en
    \texttt{UseCases/} (\texttt{StateT}, \texttt{ReaderT}, \texttt{ExceptT}). El bucle en
    tiempo real interpreta el paso de simulación, aplica la entrada del \emph{player} y
    mantiene el \emph{run} entre niveles.

  \item Cargar \emph{level definition}s desde JSON: geometría de plataformas,
    \emph{spawn point}, \emph{enemy}, ítems, \emph{exit zone} y \emph{hazard}s.
    Opcionalmente, generar etapas vía API y resolver \texttt{behaviourHint} al armar el
    catálogo del \emph{run}. \texttt{GameConfig} (gravedad, vidas, combate) vive en
    código, no en JSON.
\end{enumerate}

Como complemento propio del equipo, se publicó una versión ejecutable del juego en
itch.io para facilitar su prueba fuera del entorno de desarrollo.

\subsection{Metodología de trabajo}
\label{subsec:intro-metodologia}

El desarrollo siguió un enfoque incremental por funcionalidad: primero movimiento y física
del \emph{player}, luego plataformas y colisiones, después ítems y combate, enemigos con
comportamiento programable, y por último carga de niveles desde datos externos. Cada
iteración cerraba un ciclo completo (modelo de dominio, integración en el \emph{frame} y
prueba manual en el juego) antes de avanzar al siguiente bloque.

Los cambios entraron al repositorio mediante \emph{pull requests} revisadas entre
integrantes. En esas revisiones comprobamos que el bloque compilara, que la lógica nueva
respetara la frontera de pureza de \texttt{Domain/}, que los casos del dominio quedaran
cubiertos en el \emph{pattern matching}, y que los \emph{smoke checks} agregados
reprodujeran el comportamiento esperado. Esa code review complementó la validación
visual en pantalla.

Un pipeline de continuous integration en GitHub ejecuta en cada cambio la compilación con
tests habilitados, la suite \texttt{tasty}, la validación del paquete Cabal, la
documentación Haddock, HLint y Fourmolu. Esos chequeos automáticos detectan regresiones de
compilación, roturas en tests y desvíos de estilo antes del merge.

Durante el desarrollo también recurrimos a Cursor y Claude. La arquitectura del motor, el
modelado del dominio y las decisiones que este informe desarrolla en los capítulos
siguientes las tomamos entre ambos integrantes. Esas herramientas nos ayudaron a barajar
alternativas, a contrastar técnicas de Programación Funcional y a señalar debilidades en
borradores intermedios. En ningún caso delegamos la decisión final: cada elección de diseño
quedó argumentada y acordada por el equipo.

En la escritura del código la asistencia se limitó sobre todo a sintaxis de Haskell y a
piezas repetitivas de estructura. Todo lo que incorporamos pasó por las revisiones entre
integrantes y por el pipeline de CI: lo releímos, lo corregimos y lo validamos contra el
comportamiento esperado en pantalla o en tests.
